\documentclass[../../../include/open-logic-section]{subfiles}

\begin{document}

\olfileid{sth}{spine}{valpha}

\olsection{مرحلیاں دی $V_\alpha$ دے روپ وچ تعریف}

\olref[sth][ordinals][]{chap} وچ اسیں حُسنِ ترتیب والیاں ساختاں تے
(فون نوئمان دے) ترتیبی عدد تعریف کیتے سن۔ ایس باب وچ اسیں ایہناں نوں
سیٹاں دی درجہ بندی \emph{آپ} دی خاصیت دسن لئی ورتاں گے۔ ایس لئی یاد کرو کہ
\olref[sth][ordinals][opps]{sec} وچ اسیں جانشین تے حدّی ترتیبی عدداں دا تصور
تعریف کیتا سی۔ ہِیٹھلی تعریف وچ اسیں ایہ تصور ورتدے آں:

\begin{defn}\ollabel{defValphas}
	\begin{align*}
	V_\emptyset &\defis \emptyset\\
	V_{\ordsucc{\alpha}} &\defis \Pow{V_\alpha} & &
	\text{کسے وی ترتیبی عدد لئی }\alpha\\
	V_{\alpha} &\defis \bigcup_{\gamma < \alpha} V_\gamma & &
	\text{جدوں }\alpha\text{ حدّی ترتیبی عدد ہووے}
\end{align*}
\end{defn}
\noindent
ایہ ترتیبی عدداں اُتے \emph{لامتناہی درجیاں والی بازگشت} نال تعریف ہووے گی۔
ایس لحاظ نال سانوں ایہدا قدرتی عدداں اُتے فنکشوناں دیاں بازگشتی تعریفاں
نال ٹاکرا کرنا چاہیدا اے۔\footnote{\olref[sfr][infinite][dedekind]{sec}
وچ جمع، ضرب تے قوت دیاں تعریفاں نال ملا کے ویکھو۔} قدرتی عدداں وانگ
اسیں اک مُڈھلی صورت تے جانشین صورتاں تعریف کردے آں۔ پر ترتیبی عدداں
نال کم کردیاں سانوں \emph{حدّی} صورتاں وی دسّنیاں پیندیاں نیں۔

$V_\alpha$ والیاں ایہ تعریفاں اک اہم پڑاؤ ہون گیاں۔ اسیں غیر رسمی طور
اُتے سیٹاں دی درجہ بندی دی وجہ ایویں دسی اے کہ سیٹ \emph{مرحلیاں} وچ بندے
نیں۔ $V_\alpha$ والے سیٹ حقیقت وچ اوہی مرحلے نیں۔ پر ضروری گل ایہ اے
کہ ایہ مرحلیاں دی \emph{اندرلی} خاصیت اے۔ اسیں نظریے دا اک ممکن
\emph{نمونہ} پیش کرن دی تھاں مرحلے اپنے سیٹ نظریے \emph{اندر} تعریف
کر دتے ہون گے۔

\end{document}
